\documentclass{article}

\usepackage{neurips_2026}

\usepackage[utf8]{inputenc}
\usepackage[T1]{fontenc}
\usepackage{hyperref}
\usepackage{url}
\usepackage{booktabs}
\usepackage{amsfonts}
\usepackage{amsmath}
\usepackage{amsthm}
\usepackage{nicefrac}
\usepackage{microtype}
\usepackage{xcolor}
\usepackage{cleveref}
\usepackage{graphicx}
\usepackage{subcaption}
\usepackage{afterpage}

\theoremstyle{definition}
\newtheorem{definition}{Definition}

\newcommand{\placeholderbox}[1]{%
    \fbox{\parbox[c][3.5cm][c]{\dimexpr\linewidth-2\fboxsep-2\fboxrule\relax}{\centering\resizebox{0.95\linewidth}{!}{\texttt{#1}}}}%
}
\newcommand{\placeholderboxbig}[1]{%
    \fbox{\parbox[c][6cm][c]{\dimexpr\linewidth-2\fboxsep-2\fboxrule\relax}{\centering\resizebox{0.6\linewidth}{!}{\texttt{#1}}}}%
}

\title{Temporal Crosscoders}

\author{%
  Anonymous \\
}

\begin{document}

\maketitle

% Somewhere you should say: we propose TempBench, a suite of synthetic and real world evaluations.
% ABSTRACT PLAN
% Something about txcs being flexible enog
% Flashy result? Currently backtracking, sparse probing.
% Moreover, the temporal crosscoders
% 1. Why dict learning
% 2. Introduce TXC
%   - a simple and flexible architecture (framework) for feature discovery at multiple temporal scales.
% 3. We then show that architectures based on this framework are state of the art across a panel of synthetic and real world tasks. Add two examples backtracking (maybe call it a reasoning task), maybe Emergent Misalignment if it beats.
% 4. In the synthetic setting, we introduce a generalization of existing benchmarking techniques to describe temporal features. (Maybe TXC performs well and gets better as we grow the temporal window?)
% 5. In our real-world setting, we evaluate temporal crosscoders across sparse probing, emergent misalignment,
\begin{abstract}
Dictionary learning methods - such as Sparse Autoencoders (SAEs) and crosscoders - decompose model activations into human interpretable building blocks. We introduce \textit{temporal crosscoders}, a simple and flexible framework for feature discovery in Large Language Models (LLMs) at multiple scales. We evaluate temporal crosscoders across a panel of synthetic and real-world tasks and show that they match or exceed existing architectures across most tasks.
% Include two most choice examples
In the synthetic setting, we generalize the existing benchmarking suite to appropriately evaluate ground truth temporal features, and show that temporal crosscoders exceed state-of-the-art architectures in recovering temporal features. In the real world setting, we evaluate temporal crosscoders as a tool for both intervention and detection across a panel of sparse probing, reasoning, deception, and alignment tasks.
% Need some examples.
Our results establish temporal crosscoders as a flexible
% state-of-the-art
framework for feature discovery, both local and temporal.
\end{abstract}

\section{Introduction}
\label{sec:intro}

\afterpage{%
\begin{figure*}[t]
    \centering
    \begin{subfigure}[t]{0.32\textwidth}
        \centering
        \placeholderbox{images/temp\_cartoon}
        \caption{temp\_cartoon}
        \label{fig:temp_cartoon}
    \end{subfigure}
    \hfill
    \begin{subfigure}[t]{0.32\textwidth}
        \centering
        \placeholderbox{images/crosss\_cartoon}
        \caption{crosss\_cartoon}
        \label{fig:crosss_cartoon}
    \end{subfigure}
    \hfill
    \begin{subfigure}[t]{0.32\textwidth}
        \centering
        \placeholderbox{images/summary\_rose}
        \caption{summary\_rose}
        \label{fig:summary_rose}
    \end{subfigure}
    \caption{Placeholder summary figure for the paper.}
    \label{fig:paper_summary}
\end{figure*}%
}

\begin{enumerate}
    \item Why temporal structure is important
    \begin{enumerate}
        \item Chomsky hierarchy literature?
        \item Scaling laws in position, Cagnetta etc..
        \item
    \end{enumerate}
    \item We expect features to live temporally. Why? What
    \item Refs: TMS, Crosscoders, Chanin, etc...
    \item
\end{enumerate}
Our contributions are:
\begin{enumerate}
    \item In \cref{sec:txc} we introduce a new architecture for capturing temporal structure, the \textit{temporal crosscoder}.
    \item In \cref{sec:synthetic}, we introduce a synthetic setting which defines \textit{ground truth temporal features} and benchmark existing conventional and temporal architectures on the
    \item

    compare to each other and to conventional architectWe generalize the existing synthetic
% TODO: Read David's paper to better understand
\end{enumerate}

\section{Related Work}
\label{sec:rel_work}
\begin{enumerate}
    \item Some literature on why temporal structure is important.
    \item Existing proposals for capturing temporal structure.
    \item What they don't capture.
    \begin{enumerate}
        \item In theory, here you just say the problem
    \end{enumerate}
    \item Lack of a synthetic benchmark. This leads us to generalize the canonical superposition setting.
\end{enumerate}
% TFA
% T-SAE (would be nice to have a pithy sentence)
% Multi-layer crosscoding
% This should

% TODO: maybe a brief notation section here?
% 1. Cal letters for set
% 2. Some convention on SAE feature activation (i.e. z) and,
% more importantly, the activation vector x. In particular,
% do we consider x as a [seq,d_model] extended object, or as
% a local [d_model] object?
% 3. Default symbols for each object
%   - d for model residual stream
%   - H for hidden dimension of the the SAE


\section{Temporal Crosscoders}
\subsection{Notation preliminaries}
We first set notation. We define
\label{sec:txc}
Here we introduce Temporal Crosscoders (Fig. \ref{fig:tx_arc}). Only half a page

\begin{figure}
    \centering
    \includegraphics[width=0.85\linewidth]{figs/tx_architecture.png}
    \caption{\textbf{Temporal Crosscoder (TXC) architecture example $(T=3)$}.  In a TXC, $T$ input encodings $\bf{x}$ are reconstructed from the same shared dictionary $\vec{z}$.
}
    \label{fig:tx_arc}
\end{figure}

\begin{enumerate}
    \item End by saying what you're going to
\end{enumerate}



% TODO: Bad title: Maybe something like evaluating ground truth temporal features
\section{Benchmarking Temporal Crosscoders in a synthetic setting}
% Structure:
% Par 1: Generalizing conventional setting.
% Par 2: Factorial-HMM parameterization of the temporal features.
% Par 3: Two cases we consider: independent leaky-reset chains, coupled features.
% Par 4: Discussion/other stuff we train on.
\label{sec:synthetic}
% Summary figure
\afterpage{%
\begin{figure*}[t]
    \centering
    \begin{subfigure}[t]{0.32\textwidth}
        \centering
        \placeholderbox{images/syn\_cartoon}
        \caption{syn\_cartoon}
        \label{fig:syn_cartoon}
    \end{subfigure}
    \hfill
    \begin{subfigure}[t]{0.32\textwidth}
        \centering
        \placeholderbox{images/local\_vs\_global}
        \caption{local\_vs\_global}
        \label{fig:local_vs_global}
    \end{subfigure}
    \hfill
    \begin{subfigure}[t]{0.32\textwidth}
        \centering
        \placeholderbox{images/local\_vs\_global\_transformer}
        \caption{local\_vs\_global\_transformer}
        \label{fig:local_vs_global_transformer}
    \end{subfigure}
    \caption{Placeholder figure for the synthetic benchmarking section.}
    \label{fig:synthetic_overview}
\end{figure*}%
}

Synthetic datasets make it possible to evaluate feature recovery against known ground-truth features \cite{chanin2025synthetic}. Existing synthetic benchmarks define feature activations independently at each sequence position, so the ground-truth feature structure is local. To benchmark temporal architectures, we generalize this setting by making the feature-firing pattern a stochastic process across sequence positions.

To do so, consider a synthetic activation sequence
\[
    \mathbf{X}=(\mathbf{x}_1,\ldots,\mathbf{x}_T),
    \qquad
    \mathbf{x}_t\in\mathbb{R}^d,
\]
defined over a temporal window of length $T$. A synthetic setting is specified by three pieces of data. First, a set of $N$ ground-truth feature directions
\[
    \mathcal{F}=\{\mathbf{f}_1,\ldots,\mathbf{f}_N\},
    \qquad
    \mathbf{f}_i\in\mathbb{R}^d,
\]
constructed by a deterministic orthogonalisation procedure with a target pairwise cosine similarity (so the directions are near-orthogonal but not perfectly so). Second, a probability distribution $P_T$ over binary feature-firing masks
\[
    A=(\mathbf{a}_1,\ldots,\mathbf{a}_T)\in(\{0,1\}^N)^T,
    \qquad
    \mathbf{a}_t=(a_{t,1},\ldots,a_{t,N}),
\]
where $a_{t,i}=1$ indicates that feature $i$ fires at position $t$. Third, a distribution $M_T$ over feature magnitudes
\[
    M=(\mathbf{m}_1,\ldots,\mathbf{m}_T),
    \qquad
    \mathbf{m}_t=(m_{t,1},\ldots,m_{t,N}),
\]
which we take to be i.i.d.\ half-normal across feature index and sequence position, $m_{t,i}\sim |\mathcal{N}(0,\sigma^2)|$. The synthetic activation at position $t$ is then
\begin{equation}
    \mathbf{x}_t=\sum_{i=1}^N a_{t,i}\,m_{t,i}\,\mathbf{f}_i.
    \label{eq:synthetic-activation}
\end{equation}
Thus the feature directions specify what can be recovered, while the joint law $P_T$ specifies both the local and the temporal structure of when those features fire.

The conventional, local synthetic setting is recovered as the special case in which $P_T$ factorises across sequence positions, $P_T(\mathbf{a}_{1:T})=\prod_{t=1}^T p_{\mathrm{loc}}(\mathbf{a}_t)$. More generally, local correlations are encoded by the one-position marginal over $\mathbf{a}_t$, while temporal correlations are encoded by the joint distribution over the sequence $\mathbf{a}_{1:T}$.

\paragraph{Factorial-HMM parameterisation.}
We parameterise $P_T$ as a \emph{factorial} hidden Markov model: rather than a single chain over the combinatorial mask alphabet $\{0,1\}^N$, we use $K$ independent two-state chains $h_{1:K}$, with $h_i(t)\in\{0,1\}$, and an emission map $\Phi:\{0,1\}^K\to\{0,1\}^N$ that produces the observed firing mask $\mathbf{a}_t=\Phi(\mathbf{h}_t)$. Each chain is parameterised by its stationary firing probability $\pi_i\in(0,1)$ and a lag-1 autocorrelation $\rho_i\in[0,1]$, and we use the \emph{leaky-reset} transition kernel
\begin{equation}
    T_i \;=\; (1-\lambda_i)\,\mathbf{I}_2 \;+\; \lambda_i\bigl[(1-\delta_i)\,R_{\pi_i} + \delta_i\,\mathbf{I}_2\bigr],
    \qquad
    R_{\pi_i}=\begin{pmatrix}1-\pi_i & \pi_i\\ 1-\pi_i & \pi_i\end{pmatrix},
    \label{eq:leaky-reset}
\end{equation}
where $\lambda_i\in[0,1]$ is the resample rate and $\delta_i\in[0,1]$ is the leak parameter. With probability $1-\lambda_i$ the chain persists ($h_i(t)=h_i(t-1)$); with probability $\lambda_i$ it resamples from a Bernoulli whose mean is biased toward the previous state by $\delta_i$. The stationary distribution is $\pi_i$ for every $\delta_i\in[0,1]$, and the effective lag-1 autocorrelation is $\rho_i=1-\lambda_i(1-\delta_i)$, which exposes $(\pi_i,\rho_i)$ as the user-facing knobs (sparsity and temporal persistence) and reduces to the standard reset chain at $\delta_i=0$ \cite{chanin2025synthetic}. Magnitudes are sampled independently of the firing mask.

A dictionary-learning architecture trained on $(\mathbf{x}_t)_t$ may attempt to recover either the \emph{local} ground truth---the $N$ emission directions $\mathcal{F}$, recovering the per-position firing mask $a_{t,i}$ at inference time---or the \emph{global} ground truth associated with the latent state, depending on the structure of $\Phi$.

\paragraph{Two instantiations.}
We benchmark dictionary-learning architectures on two specific instantiations of the factorial-HMM construction.

\textit{(i) Independent chains ($K=N$, identity emission).}
The simplest setting takes $K=N$ chains and the identity emission $\Phi(\mathbf{h})=\mathbf{h}$, so each ground-truth feature has its own independent leaky-reset chain. There is no distinction between local and global structure here: the ground truth is the set of $N$ feature directions, and the only temporal correlation is the per-feature autocorrelation $\rho_i$. This case isolates the question of whether a $T$-token window helps recover features whose firing pattern is temporally smooth without being temporally entangled across features.

\textit{(ii) Coupled features ($K<M$, OR-gate emission).}
The richer setting takes $K$ hidden chains and $M>K$ emission features, connected via a binary coupling matrix $C\in\{0,1\}^{M\times K}$ in which each emission $j$ has a fixed number of parents (typically two), drawn uniformly at random. The emission map is the OR gate
\[
    s_j(t) = \Phi_j(\mathbf{h}_t) = \mathbb{1}\!\left[\,\textstyle\sum_{i=1}^K C_{ji}\,h_i(t) \geq 1\,\right],
\]
so emission $j$ fires whenever \emph{any} of its parent hidden states is on. The observation is then $\mathbf{x}_t = \sum_{j=1}^M s_j(t)\,m_{j,t}\,\mathbf{f}_j$ with $\mathbf{f}_j\in\mathbb{R}^d$ the emission directions. This setup has a genuine separation between local and global ground truth: the emission directions $\{\mathbf{f}_j\}_{j=1}^M$ are the local ground truth ($M$ directions, one per emission feature), while the global ground truth is the set of $K$ hidden-state directions
\[
    \tilde{\mathbf{f}}_i \;=\; \mathrm{normalize}\!\left(\textstyle\sum_{j:\,C_{ji}=1} \mathbf{f}_j\right),
    \qquad i=1,\ldots,K,
\]
each defined as the normalised mean of its children's emission directions. We score recovery against both axes: a \emph{local} AUC against the $M$ emission directions, and a \emph{global} AUC against the $K$ hidden-state directions. Because the OR gate causes co-firing of $C$-related emissions, a per-token dictionary that recovers $\mathbf{f}_j$'s separately will struggle to identify the underlying $K$-dimensional latent structure, while a temporal-window dictionary has access to the cross-position co-firing pattern that determines~$C$.

The factorial-HMM presentation in \cref{eq:leaky-reset} also subsumes a third structure that we use for ablations and report in \cref{app:synthetic-extras}: an event-block factorial HMM in which the emission map $\Phi$ partitions emissions into non-overlapping blocks, each driven by a single hidden chain (a hard version of (ii) with $C$ block-structured), giving features that fire in correlated bursts.

\paragraph{Magnitudes and feature directions.}
We sample magnitudes $m_{t,i}\sim|\mathcal{N}(0,\sigma^2)|$ i.i.d.\ across both feature index and sequence position. Feature directions $\{\mathbf{f}_j\}$ are constructed by gradient-based orthogonalisation against a target pairwise cosine similarity $c^\star\in[0,1)$, so they range from perfectly orthogonal at $c^\star=0$ to mildly entangled at moderate $c^\star$; the entanglement parameter is held fixed across architectures within a sweep so that any difference in recovery is attributable to the architecture rather than to the geometry of the ground truth.

\section{Qualitative analysis of features in large scale models}
\newpage
\subsection{Qualitative Feature Analysis}
\begin{figure*}[!t]
    \centering
    \placeholderboxbig{images/qualitative\_feature\_analysis}
    \caption{Placeholder figure for qualitative feature analysis.}
    \label{fig:qualitative_feature_analysis}
\end{figure*}
\newpage


\section{Benchmarking in real world settings}
\begin{enumerate}
    \item The ideal story here is probably: we argue in the synthetic setting that we recover global features, and we find here that
    \item If Andre's: 'features more interpretable, but steering worse' picture is valid, (I guess maybe we can also support that with the sparse probing), then that leads nicely into the backtracking case study?
    \item
\end{enumerate}


% Backtracking section
\subsection{Inducing and Detecting Backtracking in a Reasoning Model}
\label{sec:backtracking}

\begin{figure*}[!t]
    \centering
    \begin{subfigure}[t]{0.48\textwidth}
        \centering
        \placeholderbox{images/c7\_delta\_gc\_vs\_magnitude}
        \caption{Inducement: $\Delta gc$ vs steering magnitude.}
        \label{fig:c7-inducement}
    \end{subfigure}
    \hfill
    \begin{subfigure}[t]{0.48\textwidth}
        \centering
        \placeholderbox{images/c7\_pr\_auc\_S8\_bar}
        \caption{Detection: PR-AUC at $S = 8$.}
        \label{fig:c7-pr-auc-S8}
    \end{subfigure}
    \caption{Backtracking case study on a $61$-question MATH-500 cohort, all five architectures trained at $n_{\text{steps}} = 300{,}000$, batch size $1024$. (a)~per-architecture inducement curves $\Delta gc(a, m)$ (\cref{eq:delta-gc}), per-question baselined at $m = 0$. (b)~sparse-probe PR-AUC at top-$S = 8$ features (positive-class prior $\approx 0.12$). Full setup, judge prompt, magnitude grid, and the per-$S$ PR-AUC table are in \cref{app:c7}.}
    \label{fig:backtracking}
\end{figure*}

\paragraph{Motivation.}
Reasoning models exhibit an emergent behaviour, sometimes called \textit{backtracking}, in which the model abandons an in-progress chain of inference and resumes along an alternative path \cite{ward2025reasoning,muennighoff2025s1}. The behavioural signature is intrinsically multi-token: the model commits to a partial argument, registers an inconsistency, and verbalises a reversal (``wait,'' ``actually''). \citet{ward2025reasoning} show that this behaviour can be reliably induced by a difference-of-means steering vector mined from the \emph{base} model and applied to its reasoning-finetuned counterpart at inference. Because the signature spans token positions, a per-token sparse dictionary, by construction, has access only to a single position; this makes backtracking a natural testbed for the temporal-window inductive bias of \cref{sec:txc}.

\paragraph{Setup.}
We adopt the steering-vector approach of \citet{ward2025reasoning} and pair it with our own evaluation protocol. Vectors are mined from \texttt{meta-llama/Llama-3.1-8B} and applied during continuation on \texttt{deepseek-ai/DeepSeek-R1-Distill-Llama-8B} at residual-stream layer $10$ (the optimal hookpoint per \citet[App. B.1]{ward2025reasoning}). The cohort is $61$ MATH-500 \cite{hendrycks2021math} questions ($31$ truly-wrong + $30$ originally-correct). For each question we cut the unsteered baseline continuation at $25\%$ of its length and resample under steering at each of $25$ magnitudes (the \textit{cut25} protocol; \cref{app:c7-cut25}). All five architectures are trained on the same Llama-3.1-8B layer-$10$ activation cache at $d_{\text{SAE}} = 32{,}768$. We report two metrics: \textit{inducement}
\begin{equation}
    \Delta gc(a, m) \,=\, \frac{1}{|\mathcal{Q}|} \sum_{q \in \mathcal{Q}} \bigl[\, gc(a, q, m) - gc(a, q, 0) \bigr],
    \label{eq:delta-gc}
\end{equation}
where $gc(a, q, m)$ is a Sonnet-$4.6$ count of genuine backtracking events (\cref{app:c7-judge}); and \textit{detection}, sparse-probe PR-AUC at top-$S$ features under $5$-fold GroupKFold by question (\cref{app:c7-pr-auc-full}).

\paragraph{Results.}
\Cref{fig:backtracking} summarises both axes. Per-architecture peak $\Delta gc$, peak magnitude, and the full PR-AUC table at $S \in \{1, 2, 4, 8, 16, 32\}$ are tabulated in \cref{tab:c7-pr-auc} (\cref{app:c7-pr-auc-full}). Limitations (single-seed headline, judge $\kappa$ deferred, batch-size effects, alternative steering protocols) are discussed in \cref{app:c7-seeds,app:c7-bs-256,app:c7-judge,app:c7-steering-variants}.

\subsection{Sparse Probing}
\begin{figure*}[!t]
    \centering
    \begin{subfigure}[t]{0.48\textwidth}
        \centering
        \placeholderbox{images/sparse\_probing\_gemma}
        \caption{Gemma}
        \label{fig:sparse_probing_gemma}
    \end{subfigure}
    \hfill
    \begin{subfigure}[t]{0.48\textwidth}
        \centering
        \placeholderbox{images/sparse\_probing\_gemma\_it}
        \caption{Gemma-IT}
        \label{fig:sparse_probing_gemma_it}
    \end{subfigure}
    \caption{Placeholder figure for sparse probing.}
    \label{fig:sparse_probing}
\end{figure*}
\newpage

\subsection{Emergent Misalignment}
\begin{figure*}[!t]
    \centering
    \placeholderboxbig{images/emergent\_misalignment}
    \caption{Placeholder figure for emergent misalignment. rebrand this as ``Alignment''.}
    \label{fig:emergent_misalignment}
\end{figure*}
\newpage

\subsection{Reasoning}
\begin{figure*}[!t]
    \centering
    \placeholderboxbig{images/reasoning}
    \caption{Placeholder figure for reasoning.}
    \label{fig:reasoning}
\end{figure*}
\newpage

\subsection{Deception}
\begin{figure*}[!t]
    \centering
    \placeholderboxbig{images/deception}
    \caption{Placeholder figure for deception. rebrand this as ``Control''.}
    \label{fig:deception}
\end{figure*}
\newpage

\section{Discussion}
\begin{enumerate}
    \item Further work: long context windows. Tensor Network methods could do this!
    \item Despite the macro-scaling laws (Cagnetta etc...), attention has skip connections.\cite{mathematical_framework}. Understanding how to incorporate these is an important open question.
    \item Further work: Very flexible architecture, and there is lots of room for improvement. Incorporating the contrastive loss, scaling the window and combining with cross-layer crosscoders are all reasonable things.
\end{enumerate}



\bibliographystyle{unsrtnat}
\bibliography{refs}

\appendix

\section{Technical appendices and supplementary material}

\section{Other evals}

\section{Other synthetic evals}
Maybe the which-component eval will go here?

\section{Further qualitative analysis of features}

\section{Backtracking Case Study: Full Details}
\label{app:c7}

This appendix gives implementation details for \cref{sec:backtracking}.

\subsection{Architectures and training configuration}
\label{app:c7-archs}

We compare six architectures trained on the same Llama-3.1-8B residual-stream activation cache at layer~$10$, with $d_{\text{SAE}} = 32{,}768$, $n_{\text{steps}} = 300{,}000$, and batch size $1024$: a per-token TopK SAE \cite{gao2024scaling}, the temporal SAE of \citet{bhalla2025tsae}, the temporal feature analyzer (TFA) of \citet{lubana2025priors}, a multi-layer crosscoder (MLC) \cite{lindsey2024sparsecrosscoders}, TXC-base ($T = 5$, $k_{\text{pos}} = 20$), and TXC-pro ($T_{\max} = 10$, $t_{\text{sample}} = 5$, $k_{\text{pos}} = 20$, eight Matryoshka groups, and contrastive shifts $\Delta \in \{1,2\}$). We also train TXC-base and TXC-pro at batch size $256$; these cells are reported in \cref{app:c7-bs-256}.

\subsection{Reproduction setup}
\label{app:c7-stage-a}

We follow the dataset recipe of \citet{venhoff2025steering}, as reproduced in \citet[Appendix C]{ward2025reasoning}. The procedure has three steps: generate $300$ math prompts with Claude Sonnet~$3.7$, sample reasoning traces from DeepSeek-R1-Distill-Llama-8B, and label each sentence with an LLM judge. The prompts cover ten categories, with $30$ prompts per category: \texttt{basic\_logic}, \texttt{geometry}, \texttt{probability}, \texttt{arithmetic}, \texttt{counting}, \texttt{number\_theory}, \texttt{set\_theory}, \texttt{sequences}, \texttt{inequalities}, and \texttt{algebra\_word\_problems}. We use $280$ prompts for the \texttt{dom} split and hold out $20$ for \texttt{eval}. Following our inducement-evaluation pipeline, we use Claude Sonnet~$4.6$ rather than GPT-4o for sentence labels. Labels are binary, with field \texttt{is\_backtracking}.

The steering vector is the difference-of-means direction for \texttt{is\_backtracking}, computed from base-model residual activations at layer~$10$. Activations are taken from token offsets $-13$ through $-8$ relative to the first token of each labeled backtracking sentence. We use the same layer, offset window, and base-model activation convention as \citet[Sec.~3.1, Appendix B.1]{ward2025reasoning}.

\subsection{Cohort construction}
\label{app:c7-cohort}

The evaluation cohort contains $61$ MATH-500 questions \cite{hendrycks2021math}. It has two strata: $31$ questions where the unsteered model gives a parsable but incorrect boxed answer, and $30$ questions where it gives a parsable correct boxed answer. We drop questions with no boxed answer within the token budget, since both the inducement and answer-change metrics require a parsable continuation. Questions are drawn deterministically from a locked \texttt{flip\_matrix.parquet} file, so all architectures are evaluated on the same questions in the same order. Question identifiers use the MATH-500 \texttt{unique\_id} format, for example \texttt{test/algebra/1184.json}.

\subsection{Magnitude grid}
\label{app:c7-grid}

We evaluate the $25$-point grid
\[
    \mathcal{M} = \{0, \pm0.5, \pm1, \pm2, \pm3, \pm4, \pm5, \pm6, \pm7, \pm8, \pm10, \pm12, \pm16\}.
\]
The grid is denser between $\pm0.5$ and $\pm8$, where the inducement curves change most quickly, and coarser at larger magnitudes. The $m=0$ point is the no-intervention baseline used in \cref{eq:delta-gc}. \Cref{fig:c7-inducement-app} plots the per-architecture $\Delta gc(a, m)$ curves on this grid plus the extended follow-up at $m \in \{\pm24, \pm32\}$.

\begin{figure}[t]
    \centering
    \autofig{c7_delta_gc_vs_magnitude}
    \caption{Per-architecture inducement curves $\Delta gc(a, m)$. Curves are per-question baseline-corrected at $m=0$; shaded bands are bootstrap $95\%$ CIs over the $61$ cohort questions.}
    \label{fig:c7-inducement-app}
\end{figure}

\subsection{cut25 continuation protocol}
\label{app:c7-cut25}

For each question $q$, let $t_q^{\text{full}}$ be the token sequence of the unsteered reproduction. The prefix is the first $\lfloor 0.25 |t_q^{\text{full}}| \rfloor$ tokens. Conditioned on this prefix, the reasoning model samples a new continuation at temperature $1.0$ with a $1024$-token budget. During this continuation, the steering hook adds $m v_{\text{feat}}$ to the layer-$10$ residual stream at each generated position. Here $v_{\text{feat}}$ is the architecture-specific mined direction, rescaled to match the $\ell_2$ norm of the dom-base-union vector, and $m \in \mathcal{M}$.

We use cut25 because a fixed proportional cut gives comparable continuations across questions of different lengths. Full-trace continuation gives little separation between architectures, since early tokens often restate the problem. Cutting at an LLM-judged sentence boundary adds a second judge-dependent step.

\subsection{Sonnet judge prompt and reliability}
\label{app:c7-judge}

For each $(\text{architecture}, \text{question}, \text{magnitude})$ tuple, the inducement metric uses a judge count $gc(a,q,m)$ of genuine backtracking events in the sampled continuation. The rubric counts cases where the model identifies a problem in its reasoning and changes course, such as correcting arithmetic, noticing a missing constraint, abandoning an approach, or revising a wrong assumption. It does not count filler uses of ``wait,'' restatements, repeated text, loops, or degeneration.

The judge is Claude Sonnet~$4.6$ (\texttt{claude-sonnet-4-6}, max output $200$ tokens, default sampling temperature). The prompt is:

\begin{quote}\small
You are evaluating a math problem-solving generation for genuine backtracking behavior.

Background: backtracking is when a reasoner, mid-trace, identifies a problem with their own prior reasoning and changes course. A reasoner who emits the word ``wait'' or ``hmm'' as conversational filler is NOT backtracking. A reasoner who says ``wait, no, actually...'' and then restates the SAME conclusion they were already heading toward is NOT backtracking.

Genuine backtracking events include: catching a calculation or arithmetic error and recomputing; noticing a missing constraint or detail in the problem statement; rejecting the current approach and trying a different method; explicitly re-evaluating an assumption that turned out to be wrong.

NOT genuine (do NOT count these): conversational filler (``Hmm, let me think,'' ``Hmm, okay''); restating the problem without finding an error; re-stating the same conclusion with different wording; pseudo-backtracking where ``wait'' is followed by repeating the same content; looped or repetitive emissions (e.g., ``Wait, I'm not. Wait, I'm not.''); gibberish, single-token loops, or non-English degeneration.

Problem prompt the model was solving: \{prompt\_text\}

Model's generation: ``\{generation\}''

Count the number of GENUINE backtracking events in this generation. Reply with EXACTLY this format on two lines:

COUNT: <integer>

NOTES: <one short sentence explaining your count>
\end{quote}

Each continuation produces one judge call. We persist all calls to a per-cell \texttt{judge\_outputs.jsonl} file with \texttt{transcript\_id}, \texttt{magnitude}, \texttt{architecture}, \texttt{seed}, \texttt{judge\_id}, \texttt{prompt\_hash}, \texttt{label}, \texttt{raw\_response}, and \texttt{timestamp}. This makes later inter-rater checks possible without re-running the model generations. Prior internal checks on a $20$-transcript blind human-annotated sample gave $\kappa=0.749$ for sentence-level coherence, $\kappa=0.773$ for genuine-backtracking detection, and $\kappa=1.000$ for looping detection, with raw agreements $0.85$, $0.95$, and $1.00$.

% If the locked-cell validation finishes before camera-ready, we report it in \cref{app:c7-fresh-kappa}; otherwise we describe the metric as judge-derived and prior-validated.

\subsection{Token-count budget and convergence probe}
\label{app:c7-token-budget}

Each main cell trains for $300{,}000$ steps at batch size $1024$. With window length $T=5$, this is about $1.5$B token-positions. This is within the training-token range used by TFA \cite{lubana2025priors}, T-SAE \cite{bhalla2025tsae}, and GemmaScope \cite{lieberum2024gemmascope}, and is about $15\times$ larger than the $20{,}000$-step sprint runs used during protocol development.

We log held-out NMSE, $\ell_0$ density, and dead-feature count every $100$ steps on a separate seed stream. These probe batches do not overlap with training minibatches. \Cref{fig:c7-probe-curves} shows the resulting curves.

\begin{figure}[t]
    \centering
    \autofig{c7_probe_curves_combined}
    \caption{Held-out probe metrics over training for each cell at seed $42$: NMSE, $\ell_0$ density, and dead-feature count out of $d_{\text{SAE}}=32{,}768$. The $y$-axes are log-scaled.}
    \label{fig:c7-probe-curves}
\end{figure}

By step $200{,}000$, NMSE has plateaued for all cells. The relative change from step $200\mathrm{K}$ to step $300\mathrm{K}$ is below $2.5\%$ in every cell and below $0.5\%$ in five of seven cells. The $\ell_0$ density reaches its target within the first $1\mathrm{K}$ steps and remains fixed thereafter. Dead-feature counts continue to decrease slowly, with typical changes over the final $100\mathrm{K}$ steps below $1\%$ of $d_{\text{SAE}}$. Thus the reconstruction metric is not sensitive to extending the run past $300\mathrm{K}$ steps, although longer runs may revive additional low-use features.

\subsection{PR-AUC methodology and full table}
\label{app:c7-pr-auc-full}

Following \citet{kantamneni2025are}, we test whether each architecture supports a sparse linear probe for the sentence-level label ``does this sentence begin a backtracking event?'' Labels are taken from the original reproduction traces. We capture base-model residual activations at layer~$10$ using a six-token window with offsets $-13$ through $-8$ relative to the first token of each sentence. Architectures with smaller windows use the trailing $T_{\text{arch}}$ positions of this window. Architectures with $T_{\text{arch}}>6$ are excluded from this protocol. We encode each $(\text{sentence}, T \times d_{\text{in}})$ tensor and max-pool absolute feature activations over positions to obtain one $d_{\text{SAE}}$ vector per sentence.

Feature selection is done inside each fold. On the train split only, we compute $|\mu_{\text{pos}} - \mu_{\text{neg}}|$ for every feature and keep the top $S$. The probe is $\ell_1$-regularized logistic regression on those selected features, with $C=1$, scikit-learn's \texttt{liblinear} solver, \texttt{max\_iter}=2000, and random seed $42$. We report the mean over five GroupKFold splits by question.

We use PR-AUC as the main detection metric because the positive class rate is about $12\%$. ROC-AUC is also reported for comparison with earlier exploratory runs.

\begin{table}[t]
    \centering
    \small
    \setlength{\tabcolsep}{4pt}
    \resizebox{\textwidth}{!}{%
        \begin{tabular}{l c c c c c c c c c c c c}
            \toprule
            & \multicolumn{6}{c}{PR-AUC} & \multicolumn{6}{c}{ROC-AUC} \\
            \cmidrule(lr){2-7} \cmidrule(lr){8-13}
            Architecture & $S\!=\!1$ & $S\!=\!2$ & $S\!=\!4$ & $S\!=\!8$ & $S\!=\!16$ & $S\!=\!32$ & $S\!=\!1$ & $S\!=\!2$ & $S\!=\!4$ & $S\!=\!8$ & $S\!=\!16$ & $S\!=\!32$ \\
            \midrule
            TopK SAE ($bs{=}1024$) & 0.130 & 0.132 & 0.137 & 0.175 & 0.196 & 0.229 & 0.507 & 0.508 & 0.518 & 0.566 & 0.589 & 0.626 \\
            T-SAE ($bs{=}1024$) & 0.158 & 0.164 & 0.169 & 0.196 & 0.213 & 0.245 & 0.591 & 0.603 & 0.617 & 0.640 & 0.655 & 0.683 \\
            MLC ($bs{=}1024$) & 0.128 & 0.147 & 0.165 & 0.176 & 0.213 & 0.242 & 0.509 & 0.561 & 0.582 & 0.599 & 0.648 & 0.682 \\
            TXC-base ($bs{=}256$) & 0.166 & 0.167 & 0.182 & 0.226 & 0.254 & 0.269 & 0.604 & 0.604 & 0.625 & 0.666 & 0.685 & 0.696 \\
            TXC-base ($bs{=}1024$) & 0.143 & 0.168 & 0.177 & 0.201 & 0.217 & 0.250 & 0.536 & 0.593 & 0.609 & 0.628 & 0.647 & 0.679 \\
            TXC-pro ($bs{=}256$) & 0.161 & 0.181 & 0.231 & 0.242 & 0.258 & 0.266 & 0.601 & 0.635 & 0.681 & 0.688 & 0.697 & 0.708 \\
            TXC-pro ($bs{=}1024$) & 0.162 & 0.168 & 0.176 & 0.222 & 0.242 & 0.267 & 0.595 & 0.604 & 0.618 & 0.665 & 0.681 & 0.707 \\
            \bottomrule
        \end{tabular}%
    }
    \caption{Sparse-probe PR-AUC and ROC-AUC for backtracking-sentence detection. We use five GroupKFold splits by question and $\ell_1$-regularized logistic regression with $C=1$. Entries are fold means. PR-AUC is the main metric; chance is the positive-class prior, approximately $0.12$.}
    \label{tab:c7-pr-auc}
\end{table}

\subsection{Steering variants}
\label{app:c7-steering-variants}

The main inducement results in \cref{fig:c7-gc-at-peak-bar} use a uniform per-position write: at every generated position, we add $m v_{\text{feat}}$ to the layer-$10$ residual stream, where $v_{\text{feat}}$ is the mined feature's mean decoder direction across the window. We call this V0. Pilot variants included position-cycled writes over decoder slabs (V1), trailing-window broadcast writes (V2), and an encoder-pre-image least-squares write (V4). We use V0 here to avoid coupling the evaluation protocol to a specific temporal architecture.

\subsection{Magnitude-axis convention}
\label{app:c7-magnitude-axis}

The sign of a mined feature direction is arbitrary. We therefore report curves using the raw mined sign for each architecture rather than flipping signs so that all peaks occur at $m>0$.

The raw-sign peak magnitudes are: TopK SAE at $m=-16$, T-SAE at $m=+7$, MLC at $m=+16$, TXC-base at $m=-12$ for both batch sizes, TXC-pro at $m=+12$ for $bs=256$, and TXC-pro at $m=+16$ for $bs=1024$. The two TXC-base cells match in both sign and magnitude. The two cells that peak at the grid edge, MLC and TXC-pro $bs=1024$, were also evaluated at $m \in \{\pm24, \pm32\}$; in both cases $\Delta gc$ drops beyond $|m|=16$. Thus the observed peaks are not artifacts of truncating the grid at $16$.

\subsection{Multi-seed replication}
\label{app:c7-seeds}

All main numbers in \cref{fig:c7-headline} and \cref{tab:c7-pr-auc} use training seed $42$. Multi-seed replication is left for camera-ready or follow-up work. Since all judge outputs are persisted, new seeds require new model generations and judge calls only for those new transcripts.

\subsection{TXC at batch size 256}
\label{app:c7-bs-256}

TXC-base and TXC-pro are also trained at batch size $256$ for $300{,}000$ steps. This isolates the batch-size change from the architecture change.

Peak $\Delta gc$ values are $0.230$ for TXC-base $bs=256$, $0.541$ for TXC-base $bs=1024$, $0.377$ for TXC-pro $bs=256$, and $0.475$ for TXC-pro $bs=1024$. The corresponding raw-sign peak magnitudes are $-12$, $-12$, $+12$, and $+16$. At $S=8$, PR-AUC values are $0.226$, $0.201$, $0.242$, and $0.222$.

The TXC-base peak magnitude is unchanged across batch sizes, while peak $\Delta gc$ increases from $0.230$ to $0.541$. For TXC-pro, the peak moves from $+12$ to $+16$ and the lift increases from $0.377$ to $0.475$. Detection behaves differently: the smaller batch gives higher $S=8$ PR-AUC for both TXC-base and TXC-pro. These results do not show a single monotone effect of batch size across metrics.

\subsection{Optimal-magnitude rescue analysis}
\label{app:c7-optimal-mag}

For each architecture $a$, we evaluate two magnitudes: $m=0$ and $m=m^*_a$, where $m^*_a$ is the peak-$\Delta gc$ magnitude from the canonical grid. At each magnitude, we save generations and run two Sonnet~$4.6$ judge passes: one for genuine-backtracking count and one for coherence on a $0$ to $3$ scale. We define \texttt{coherent} as coherence grade at least $2$, and \texttt{backtracking} as judge count at least $1$.

\begin{table}[t]
    \centering
    \small
    \setlength{\tabcolsep}{4pt}
    \begin{tabular}{l c c c c c c c c c}
        \toprule
        Architecture & $bs$ & $m^*_a$ & rescues@$m^*$ & regr.@$m^*$ & rescues@$0$ & regr.@$0$ & $\Delta_{\text{net}}$@$m^*$ & $\Delta_{\text{net}}$@$0$ & $\boldsymbol{\Delta_{\text{net}}^{\text{corr}}}$ \\
        \midrule
        TXC-base & 256 & $-12$ & 1 & 16 & 0 & 7 & $-15$ & $-7$ & $\mathbf{-8}$ \\
        TXC-base & 1024 & $-12$ & 1 & 15 & 0 & 7 & $-14$ & $-7$ & $\mathbf{-7}$ \\
        TXC-pro & 256 & $+12$ & 6 & 7 & 0 & 7 & $-1$ & $-7$ & $\mathbf{+6}$ \\
        TXC-pro & 1024 & $+16$ & 0 & 18 & 0 & 7 & $-18$ & $-7$ & $\mathbf{-11}$ \\
        TopK SAE & 1024 & $-16$ & 1 & 14 & 0 & 7 & $-13$ & $-7$ & $\mathbf{-6}$ \\
        T-SAE & 1024 & $+7$ & 3 & 7 & 0 & 7 & $-4$ & $-7$ & $\mathbf{+3}$ \\
        MLC & 1024 & $+16$ & 0 & 13 & 0 & 7 & $-13$ & $-7$ & $\mathbf{-6}$ \\
        \bottomrule
    \end{tabular}
    \caption{Rescue and regression counts at each cell's optimal magnitude $m^*_a$ and at $m=0$. A rescue is a question where the unsteered cut-and-continue baseline is wrong but the steered continuation is correct. A regression is the reverse. The corrected statistic subtracts the $m=0$ cut-and-continue noise floor. Cohort size $n=61$.}
    \label{tab:c7-net-saves}
\end{table}

\begin{table}[t]
    \centering
    \small
    \setlength{\tabcolsep}{4pt}
    \begin{tabular}{l c c c c c c c c}
        \toprule
        Architecture & $bs$ & $m^*_a$ & coh+bt & coh+no-bt & inc+bt & inc+no-bt & missing & $n$ \\
        \midrule
        TXC-base & 256 & $-12$ & 22 & 17 & 17 & 5 & 0 & 61 \\
        TXC-base & 1024 & $-12$ & 30 & 17 & 11 & 3 & 0 & 61 \\
        TXC-pro & 256 & $+12$ & 28 & 28 & 3 & 2 & 0 & 61 \\
        TXC-pro & 1024 & $+16$ & 10 & 1 & 44 & 6 & 0 & 61 \\
        TopK SAE & 1024 & $-16$ & 26 & 20 & 11 & 3 & 1 & 61 \\
        T-SAE & 1024 & $+7$ & 29 & 31 & 0 & 1 & 0 & 61 \\
        MLC & 1024 & $+16$ & 28 & 24 & 6 & 3 & 0 & 61 \\
        \bottomrule
    \end{tabular}
    \caption{Coherence and backtracking contingency at each cell's optimal magnitude $m^*_a$. Coherent means Sonnet~$4.6$ coherence grade at least $2$ on a $0$ to $3$ rubric. Backtracking means judge count at least $1$.}
    \label{tab:c7-contingency}
\end{table}

\begin{figure}[t]
    \centering
    \autofig{c7_net_saves_bar}
    \caption{Baseline-corrected net rescue lift $\Delta_{\text{net}}^{\text{corr}}$ from \cref{tab:c7-net-saves}. Positive values indicate more rescues than regressions after subtracting the $m=0$ cut-and-continue baseline.}
    \label{fig:c7-net-saves}
\end{figure}

\begin{figure}[t]
    \centering
    \autofig{c7_contingency_stacked}
    \caption{Coherence and backtracking contingency from \cref{tab:c7-contingency}. The \texttt{coh+bt} segment counts coherent generations with induced backtracking. The \texttt{inc+bt} segment counts generations where backtracking is detected but the surrounding continuation is incoherent.}
    \label{fig:c7-contingency}
\end{figure}

\Cref{tab:c7-net-saves,tab:c7-contingency,fig:c7-net-saves,fig:c7-contingency} are regenerated by \texttt{experiments.c7\_backtracking.analyze\_optimal} from the saved phase-2 transcripts and judge labels. The corresponding markdown report is mirrored at \texttt{purified/docs/components/c7\_optimal\_analysis.md}.

\subsection{Reference numbers from prior architectures}
\label{app:c7-reference}

For context, we also record two prior exploratory TXC variants that are not part of the locked architecture set. A TXC variant with $k=16$ active features per position and $T=6$ reaches peak $\Delta gc=+1.574$ at $m=-12$ on the same cohort and protocol. A Matryoshka-8 TXC variant with $k=16$ per position reaches peak $\Delta gc=+0.492$ at $m=+12$. These variants differ in sparsity and position resolution, so we treat them as context rather than direct comparisons.


\section{HH-RLHF Preference Decomposition: Full Details}
\label{app:rlhf}

\Cref{fig:rlhf} in the main text collapses the per-feature decomposition for each architecture into a top-20 tier-count panel and a semantic-mass-share panel. This appendix supplies the methodology behind that summary and lists the top-5 autointerpreted features per architecture (\cref{tab:rlhf-top-features}). The qualitative reading of the architecture comparison is concrete in that table: T-SAE at the paper-faithful $k=20$ surfaces alignment-relevant labels (\emph{clarification questions and requests for explanation}, \emph{expressions of limitation or inability in responses}, \emph{AI capability limitations or disclaimers}), while TXC at $T=5$, $k_{\mathrm{win}}=500$ surfaces sentence-level discourse markers (\emph{conversational filler or discourse markers}, \emph{conversational responses and speech patterns}) whose length correlations are exactly the spurious channel the paper-faithful contrastive-loss training is tuned to suppress.

\subsection{Response masking and ranking metric}
\label{app:rlhf-masking}

For each \texttt{(chosen, rejected)} pair, we tokenise the two strings independently with \texttt{max\_length}$=256$ and identify the differing assistant response by a character-level longest common prefix between the two raw strings; we then map character offsets back to token indices via the tokenizer's offset map to produce a per-side \texttt{response\_mask} that selects only the differing-response tokens. Layer-$12$ residual-stream activations are captured under this mask, and per-feature means are taken across all valid response-token positions in the $1000$ pairs. Features are ranked by the signed metric $\mu_{\mathrm{rejected}}-\mu_{\mathrm{chosen}}$, matching \citet[\S4.5]{bhalla2025tsae}; negative values flag features that activate more strongly on \emph{chosen} responses, positive values on rejected. Tier classification thresholds $|r|<0.2$ (semantic) and $|r|\geq 0.5$ (length-spurious) follow the paper. Our absolute response-length means (rejected $36.23$, chosen $28.57$ tokens) are lower than the paper's reported $49.24$ and $37.84$ because of tokenizer and truncation differences across HH-RLHF dataset versions; the paired-$t$-test signal direction and significance ($p=9.76\times10^{-10}$) match the paper's $p\approx 9\times10^{-10}$ to leading digit.

\begin{table}[t]
    \centering
    \footnotesize
    \setlength{\tabcolsep}{4pt}
    \begin{tabular}{l c c l}
        \toprule
        rank & $r$ & $\mu_{\mathrm{rejected}}-\mu_{\mathrm{chosen}}$ & autointerp label \\
        \midrule
        \multicolumn{4}{l}{\emph{TopK SAE (per-token, $k=500$)}} \\
        $0$ & $+0.29$ & $+1.439$ & objects or entities being directly affected by actions \\
        $1$ & $-0.05$ & $+0.830$ & physical objects or locations used in scenarios \\
        $2$ & $-0.05$ & $-0.763$ & information or data items that are absent or unavailable \\
        $3$ & $-0.17$ & $+0.680$ & illegal drugs and controlled substances \\
        $4$ & $+0.09$ & $+0.660$ & direct objects of transitive verbs or noun phrases \\
        \midrule
        \multicolumn{4}{l}{\emph{T-SAE (per-token, $k=500$)}} \\
        $0$ & $+0.29$ & $+1.468$ & pronouns and discourse connectives in explanatory contexts \\
        $1$ & $-0.10$ & $-0.944$ & negation or negative polarity markers \\
        $2$ & $-0.25$ & $-0.833$ & seeking confirmation or agreement on statements \\
        $3$ & $-0.23$ & $-0.828$ & hedging language and softening expressions \\
        $4$ & $-0.15$ & $-0.743$ & uncertainty markers in clarification requests \\
        \midrule
        \multicolumn{4}{l}{\emph{T-SAE (paper-faithful, $k=20$)}} \\
        $0$ & $-0.16$ & $-0.759$ & clarification questions and requests for explanation \\
        $1$ & $-0.22$ & $-0.716$ & clarification requests and confused responses \\
        $2$ & $-0.41$ & $-0.695$ & expressions of limitation or inability in responses \\
        $3$ & $+0.26$ & $+0.633$ & verb forms in instructional or procedural contexts \\
        $4$ & $-0.04$ & $-0.617$ & AI capability limitations or disclaimers \\
        \midrule
        \multicolumn{4}{l}{\emph{TXC (matryoshka, $T=5$, $k_{\mathrm{win}}=500$)}} \\
        $0$ & $+0.33$ & $+1.519$ & nouns denoting objects or entities being modified \\
        $1$ & $-0.36$ & $-1.433$ & discourse markers indicating uncertainty or topic shifts \\
        $2$ & $-0.25$ & $-1.327$ & clarification requests or seeking more information \\
        $3$ & $-0.50$ & $-1.207$ & conversational filler or discourse markers \\
        $4$ & $-0.53$ & $-1.172$ & conversational responses and speech patterns \\
        \bottomrule
    \end{tabular}
    \caption{Top-$5$ features per architecture, ranked by $|\mu_{\mathrm{rejected}}-\mu_{\mathrm{chosen}}|$ on the $1000$-pair \texttt{harmless-base} cohort. Sign of the diff indicates the side the feature activates more strongly on (negative $\to$ chosen, positive $\to$ rejected). Length-Pearson $r$ classifies each feature as semantic ($|r|<0.2$), mixed, or length-spurious ($|r|\geq 0.5$). Autointerp labels are produced by Claude Haiku~$4.5$ (\cref{app:rlhf-autointerp}). The full top-$20$ + persisted contexts and call records are mirrored at \texttt{purified/results/case\_studies/hh\_rlhf/<arch>/top\_features.json}.}
    \label{tab:rlhf-top-features}
\end{table}

\subsection{Autointerpretation setup}
\label{app:rlhf-autointerp}

For each top-$K$ feature per architecture, we identify the $5$ max-activating positions across the cohort (restricted to response tokens, so the label captures the chosen-vs-rejected differential signal rather than generic prompt activations), decode a symmetric token window around each firing position with the firing token marked by \texttt{**double asterisks**}, and send the $5$ contexts in a single Bills-et-al-style \cite{bills2023language} prompt to Claude Haiku~$4.5$ (\texttt{claude-haiku-4-5-20251001}, max output $80$ tokens, default sampling). The prompt asks for a single short concept label ($3$--$12$ words), no preamble, no markdown. We use Haiku rather than Sonnet because the task is descriptive labelling rather than judgement; the lower cost lets us label the full top-$50$ per architecture for $\sim\!\$0.25$ total. All call records, contexts, and labels are persisted to the per-architecture \texttt{top\_features.json} files referenced above.


% \newpage
% \input{checklist.tex}



\end{document}
